\section{Background: Where EM Lives}
\label{sec:background}

This section fixes the conceptual object studied in the rest of the paper. An
implicit-EM site is not any hidden layer with a softmax nearby, and it is not an
entire deep network. It is a local exponentiation-and-normalization interface at
which an objective sees distances, energies, or logits and therefore produces a
responsibility-shaped gradient. The theory section asks when this local
structure survives the pullback to parameters.

\subsection{Responsibilities as Gradients}

For an objective with log-sum-exp structure over distances,
\[
    L_{\mathrm{LSE}}(d) = -\log \sum_{j=1}^K \exp(-d_j),
    \qquad
    r = \operatorname{softmax}(-d),
\]
the derivative with respect to each distance coordinate is
\[
    \frac{\partial L_{\mathrm{LSE}}}{\partial d_j} = r_j.
\]
The gradient is therefore a posterior responsibility vector: it is
nonnegative, sums to one, and assigns each component a fractional share of the
update. This is the sense in which gradient descent on an LSE distance
objective performs an implicit E-step. The responsibilities are not introduced
as auxiliary variables; they appear as the gradient that the optimizer receives.

This identity is algebraic. It does not require a particular parameterization,
architecture, or optimizer. Its consequences depend on what happens next: the
responsibility vector may reach private component parameters directly, or it
may be mixed with other gradients and pulled through learned maps before it
updates earlier weights.

\subsection{Cross-Entropy Is a Label-Clamped Responsibility Interface}

Supervised cross-entropy has the same exponentiation-and-normalization
interface, but with the observed label clamping one component. For logits
$h\in\mathbb{R}^C$,
\[
    L_{\mathrm{CE}}(h,y)
    =
    -h_y + \log \sum_{c=1}^C \exp(h_c),
    \qquad
    p = \operatorname{softmax}(h),
\]
and
\[
    \frac{\partial L_{\mathrm{CE}}}{\partial h_c}
    =
    p_c - \mathbf{1}\{c=y\}.
\]
The output softmax is thus an EM-like interface whose posterior is clamped by
the label: probability mass is removed from the true class and assigned to the
competing classes according to the model posterior.

This observation is useful but also dangerous. It is correct to say that the
output loss supplies responsibility-shaped geometry at the output logits. It is
not correct to conclude that every earlier hidden layer is itself doing EM. The
gradient that reaches an earlier layer has already passed through the head
Jacobian and any intervening nonlinearities. It need not be nonnegative, need
not sum to one, and need not remain aligned with hidden components.

\subsection{Local EM Sites Inside a Network}

The supervised architecture studied here deliberately creates an intermediate
site. The first map produces component distances $d_1,\ldots,d_K$. The
NegLogSoftmin transform
\[
    \operatorname{NLS}(d)_j
    =
    d_j + \log\sum_k \exp(-d_k)
\]
calibrates these distances so that
\[
    \exp(-\operatorname{NLS}(d)_j)=r_j.
\]
This gives the hidden layer an explicit distance-to-responsibility interface
rather than relying on the output loss to impose structure indirectly.

The distinction between an output interface and an intermediate site is central.
At the output, the supervised objective owns the softmax. At the hidden layer,
the local site exists only because we add a local auxiliary objective. Without
such an objective, the hidden coordinates may be useful features, but there is
no reason for their updates to be posterior responsibilities over hidden
components.

\subsection{Volume Control}

An LSE competition term by itself does not prevent the standard degeneracies of
mixture models. Components can collapse, become redundant, or receive
essentially no responsibility. The decoder-free encoder study addressed these
failures with volume control: a variance term prevents dead or collapsed
components, and a decorrelation term prevents all components from moving
together.

In this paper the local auxiliary objective is
\[
    L_{\mathrm{aux}}
    =
    L_{\mathrm{LSE}} + L_{\mathrm{var}} + L_{\mathrm{tc}},
\]
with the supervised objective
\[
    L_{\mathrm{total}}
    =
    L_{\mathrm{CE}}+\lambda L_{\mathrm{aux}}.
\]
The variance term is the anti-collapse part of the local mixture geometry. The
decorrelation or total-correlation term is the anti-redundancy part. Together
they play the role of a volume barrier: the components are encouraged to occupy
distinct directions rather than concentrate in the same small region.

The first empirical question is whether these local repairs still work inside a
supervised network. They do: Results I shows that volume control eliminates
dead units and reduces redundancy at the intermediate site. The second question
is harder: whether the optimizer conditioning observed in the single-layer
implicit-EM model transfers through the supervised head. Results II and III show
that it does not, except when the local EM path is isolated or made dominant.

\subsection{Locality Versus Inheritance}

The main conceptual error this paper rules out is inheritance. The
gradient-responsibility identity guarantees a local structure at the coordinates
seen by an LSE or softmax objective. It does not guarantee that the whole
network inherits EM's conditioning. Backpropagation can transport the gradient,
but transport is not preservation. Dense learned maps can mix signs and
components; objective sums can add unrelated gradients; gates can turn small
responsibilities into exactly zero updates.

Thus there are two different claims one might make about a hidden layer. A
weak, correct claim is that loss geometry at downstream softmax sites shapes
hidden representations through backpropagation. A strong, generally false claim
is that the hidden layer itself receives a posterior responsibility update. The
rest of the paper draws the boundary between these claims.

\subsection{Two Failure Classes}

The experiments are organized around two failure classes. At-site failures
occur at a local EM interface: dead components, collapsed variance, redundant
features, and diffuse or poorly calibrated responsibilities. These are the
failures volume control is designed to repair.

Structural failures occur when a responsibility gradient must travel through a
shared learned corridor before reaching the parameters it is supposed to
condition. In the supervised model, this corridor is the component-to-class map
$W_2$. The CE gradient reaching $W_1$ is pulled through $W_2$ and combined with
the local auxiliary gradient. When this foreign path dominates, the update at
$W_1$ no longer has the private, per-component structure of an EM M-step.

This split sets up the theory. The next section proves that private LSE sites
inherit a parameter-independent curvature bound, while composition through a
learned dense map forfeits that guarantee. The experiments then show the same
boundary from both sides: local volume control transfers, but optimizer
conditioning does not survive ordinary composition.
